\section{Coordination}
\label{sec:coord}

\paragraph{The \texttt{asa-team-v1} protocol.} Both agents speak a small typed
protocol: every message is an envelope (sender, type, payload, timestamp) with
types \texttt{hello}, \texttt{position}, \texttt{intention}, \texttt{claim},
\texttt{mission-update}, \texttt{request-help}, \texttt{handover} and
\texttt{ack}. Agents discover each other with a \texttt{shout} handshake; in a
two-agent run, A and B start from one \texttt{.env} and pass each other's name
via \texttt{--teammate} (a shared \texttt{.env} cannot hold both names), and
discovery is logged as \texttt{teammate\_discovered}. Validated live: A and B
pair and exchange position heartbeats ($\sim$1/s) bidirectionally.

\paragraph{What coordination buys.} Three things: \emph{belief sharing} (the
position heartbeat), \emph{target deconfliction} (on committing to a pickup an
agent broadcasts a \texttt{claim}, so the teammate does not chase the same
parcel), and \emph{mission propagation} (Agent~B interprets a mission and
forwards the resulting constraints to Agent~A). We distinguish \texttt{say}
(fire-and-forget broadcast) from \texttt{ask}: \texttt{ask} is reserved for true
synchronization points because the server imposes a 1\,s reply timeout, so it is
used sparingly rather than for routine state sharing.

\paragraph{Handover mission (one-pickup--another-deliver, 26c2\_8).} Roles are
explicit and derived from runtime configuration: A is the \emph{picker}, B the
\emph{deliverer}. Rather than negotiate a meeting point, both agents compute the
same \emph{map-derived deterministic rendezvous} (a non-delivery tile one step
from a delivery, tie-broken to centre then key), so they agree with zero
messages. The handover belief is \emph{coordinate-first}: the drop is located by
its tile, and the parcel id is only a hint, because acks omit ids and ids do not
survive a drop-and-repickup (Section~\ref{sec:bdi}). The picker's
\texttt{HandoverDeposit} plan carries to the rendezvous, drops, \emph{steps off
to free the tile}, and signals the drop only after actually vacating it
(otherwise \texttt{handover-exit-blocked}, no signal --- never send the
deliverer onto a tile the picker still occupies); the strategy also bars the
picker from self-delivering or re-grabbing its own drop. The deliverer's
\texttt{HandoverCollect} fetches the drop by coordinates and the normal
\texttt{DeliverCarried} plan delivers it. Validated live end-to-end on 26c2\_8:
the god mission agent awarded the cross-agent bonus repeatedly across successive
handover cycles.

\paragraph{Go-to-and-wait (26c2\_10).} The official repository ships no concrete
mission agent for this scenario (its \texttt{start.js} references a placeholder),
so we validated the documented semantics with a local god-observer fixture
(\texttt{scripts/fixtures/goto-wait-mission.js}), clearly a test harness rather
than a result on course infrastructure. Because the target centre may be a wall,
each agent goes to the nearest \emph{reachable} tile within the radius
(preferring not the teammate's tile, so the two do not contend), then waits for
the teammate via position heartbeats and holds together; on timeout the plan
aborts instead of falsely completing, keeping the mission active to retry.
Validated live: A\,$\to$\,(16,5), B\,$\to$\,(19,2), both within distance~3 of
(19,5); the fixture awarded the 500\,pt bonus. The first run was interpreted via
the deterministic fallback (the LLM timed out); a later VPN run confirmed a clean
\texttt{source=llm} interpretation (\texttt{llama-3.3-70b}, 2545\,ms: kind
\texttt{go\_to}, target (19,5), tolerance~3, hold-at-target) and the same bonus.
